\chapter{Narzędzia}
\label{chap:tools}

\section{System kontroli wersji (Michał Pawiłojć)}

W trakcie realizacji projektu \textit{UniGuesser} kluczową rolę odegrało zastosowanie systemu kontroli wersji Git oraz platformy GitHub. Ich wykorzystanie umożliwiło efektywne zarządzanie kodem źródłowym, koordynację pracy zespołu oraz utrzymanie przejrzystości i historii zmian w projekcie.

Git to rozproszony system kontroli wersji VCS (ang. \textit{Version Control System}), stworzony przez Linusa Torvaldsa w 2005 roku. Jego podstawowym zadaniem jest śledzenie zmian w plikach projektu, umożliwienie pracy wielu programistów nad tym samym kodem oraz szybkie cofanie się do wcześniejszych wersji w razie potrzeby. Git opiera się na repozytorium lokalnym i zdalnym, co zapewnia dużą elastyczność pracy – możliwa jest nawet praca offline z późniejszą synchronizacją.

Do podstawowych operacji należą:

\begin{itemize}
    \item \texttt{git init} – inicjalizacja nowego repozytorium,
    \item \texttt{git add} – dodanie zmian do indeksu (staging area),
    \item \texttt{git commit} – zapisanie zmian w historii projektu,
    \item \texttt{git branch} – zarządzanie gałęziami (np. tworzenie odgałęzień na nowe funkcjonalności),
    \item \texttt{git merge} – scalanie gałęzi,
    \item \texttt{git log} – przegląd historii commitów.
\end{itemize}

Stosowanie Gita pozwala na łatwe śledzenie błędów, eksperymentowanie z nowymi funkcjonalnościami bez wpływu na główny kod oraz pracę równoległą wielu osób. W przypadku zmian powodujących błąd powrót do działającej wersji jest praktycznie bezkosztowy.

Natomiast GitHub to jedna z najpopularniejszych platform umożliwiających przechowywanie zdalnych repozytoriów Git. Oferuje dodatkowe funkcje, takie jak:

\begin{itemize}
    \item pull requesty (prośby o połączenie zmian z główną gałęzią projektu),
    \item code review (wzajemna kontrola kodu w zespole),
    \item zarządzanie zgłoszeniami (issues),
    \item integracja z narzędziami CI/CD (ang. \textit{Continuous Integration/Continuous Delivery lub Deployment}) jak GitHub Actions.
\end{itemize}

W projekcie \textit{UniGuesser} GitHub pełnił rolę centralnego punktu współpracy, umożliwiając m.in. zgłaszanie błędów, przeglądanie zmian oraz automatyczne uruchamianie testów jednostkowych i wdrożeń (CI/CD).

Repozytorium projektu zostało zorganizowane w sposób modularny – z podziałem na frontend, backend i konfigurację Docker. Każdy członek zespołu pracował na oddzielnych branchach, a zmiany były scalane po przejściu procesu \textit{code review}. Takie podejście zapewniło stabilność projektu i ułatwiło jego rozwój.

\section{Konteneryzacja - Docker (Michał Pawiłojć)}

Współczesne aplikacje webowe, takie jak \textit{UniGuesser}, składają się z wielu komponentów - frontend, backend, baza danych - muszą one współdziałać w spójnym środowisku. Tradycyjna metoda instalowania zależności lokalnie prowadzi do problemów ze zgodnością, trudnością odtworzenia środowiska oraz błędami zależnymi od platformy. Rozwiązaniem tego problemu jest konteneryzacja - technika, która pozwala zbudować przenośne, spójne środowisko dla każdego komponentu aplikacji.

\subsection{Konteneryzacja}

Konteneryzacja to proces pakowania aplikacji wraz z całym jej środowiskiem uruchomieniowym (zależnościami, bibliotekami, konfiguracją, itd.) do jednego, samodzielnego pakietu zwanego \textbf{kontenerem}. Taki kontener może być uruchomiony na dowolnym systemie operacyjnym, który obsługuje kontenery, bez potrzeby rekonfiguracji lub instalacji dodatkowego oprogramowania.

Kontenery są lekką alternatywą dla maszyn wirtualnych i w odróżnieniu od nich nie zawierają pełnego systemu operacyjnego, a współdzielą jądro systemu hosta. Dzięki temu zużywają mniej zasobów, szybciej się uruchamiają i łatwiej się nimi zarządza.

Dzięki konteneryzacji możliwe jest podejście \textit{„build once, run anywhere”} - czyli raz zbudowany obraz aplikacji może działać tak samo na komputerze deweloperskim, serwerze testowym i w środowisku produkcyjnym.

\subsection{Docker}

Docker jest obecnie jednym z najpowszechniej stosowanych narzędzi do konteneryzacji aplikacji, umożliwiającym uruchamianie oprogramowania w odizolowanych, lekkich środowiskach zwanych kontenerami. Każdy kontener posiada własną warstwę systemową oraz zestaw zależności niezbędnych do działania aplikacji, co gwarantuje powtarzalność i niezależność od konfiguracji systemu gospodarza. Docker pozwala na definiowanie własnych obrazów aplikacji za pomocą plików \texttt{Dockerfile}, w których można określić wymagane biblioteki, sposób budowania projektu, zmienne środowiskowe oraz komendy uruchomieniowe. Dzięki temu proces tworzenia środowiska uruchomieniowego staje się w pełni automatyzowalny i możliwy do odtworzenia na dowolnym urządzeniu.

Istotnym elementem ekosystemu narzędzia jest Docker Engine, odpowiedzialny za budowanie i zarządzanie kontenerami, oraz Docker Hub - publiczny rejestr obrazów, ułatwiający dystrybucję przygotowanych konfiguracji. W szerszym ujęciu Docker stanowi fundament dla nowoczesnych praktyk DevOps, umożliwiając szybkie wdrażanie aplikacji, ograniczenie problemów związanych z różnymi środowiskami developerskimi oraz poprawę skalowalności. Zarządzanie kontenerami może odbywać się zarówno lokalnie, jak i w środowiskach chmurowych, co otwiera możliwość integracji z systemami orkiestracji, takimi jak Kubernetes, gdy zajdzie potrzeba rozwinięcia projektu na większą skalę.

\subsection{Docker Compose}

Projekt \textit{UniGuesser} składa się z kilku współdziałających usług:

\begin{itemize}
  \item Frontend (React),
  \item Backend (.NET),
  \item Baza danych (PostgreSQL).
\end{itemize}

Do zarządzania tymi usługami wykorzystano \textbf{Docker Compose}. Umożliwia ono deklaratywne opisanie całego środowiska w jednym pliku YAML	. W ten sposób tworzona jest wewnętrzna sieć, w której poszczególne usługi są w stanie się komunikować.

Całe środowisko można uruchomić jedną komendą:

\begin{lstlisting}[language=bash]
docker-compose up --build
\end{lstlisting}

\subsection{Korzyści wynikające z konteneryzacji}

Użycie Dockera i Composa przyniosło wiele korzyści:

\begin{itemize}
  \item \textbf{Powtarzalność środowiska} — każdy członek zespołu uruchamiał projekt dokładnie w taki sam sposób.
  \item \textbf{Izolacja komponentów} — każda usługa działała w swoim kontenerze.
  \item \textbf{Łatwe wdrażanie} — obrazy Dockera działały tak samo lokalnie, jak i na produkcji.
  \item \textbf{Integracja z CI/CD} — obrazy mogły być budowane i testowane w GitHub Actions.
\end{itemize}

\subsection{Ograniczenia}

Choć konteneryzacja znacząco uprościła zarządzanie środowiskiem, zostały również napotkane pewne trudności:

\begin{itemize}
  \item \textbf{Krzywa uczenia} — wymagana była znajomość Dockera i Composa.
  \item \textbf{Wydajność na Windowsie} — niektóre operacje były wolniejsze w WSL2.
  \item \textbf{Debugowanie} — błędy konfiguracyjne były trudniejsze do wychwycenia.
\end{itemize}

Pomimo tych wyzwań, konteneryzacja odegrała kluczową rolę w sukcesie projektu.

\section{Zarządzanie projektem - YouTrack (Karina Wołoszyn)}

Do zarządzania projektem wykorzystany został Youtrack, który jest narzędziem typu ITS (ang. \textit{Issue Tracker System}) stworzonym przez firmę JetBrains. Opiera się on na technologii Ajax i wspiera szeroki zakres metodyk zarządzania projektem, w tym podejść zwinnych, które również zostało wykonane do realizacji tego projektu.

System Youtrack umożliwia m.in. planowanie i organizacje zadań, przypisywanie ich do konkretnych członków, ustalanie priorytetów, śledzenia postępów prac, dodawanie tagów, załączników czy tworzenie raportów. Dzięki rozbudowanym możliwościom filtrowania, szacowania i określania ilości spędzonego czasu, personalizacji widoków możliwe było bieżące monitorowanie i reagowanie na pojawiające się ewentualne problemy. 

W ramach projektu przyjęto podejście oparte na zmodyfikowanej wersji metodyki Kanban, dostosowanej do potrzeb zespołu i charakteru realizowanego zadania. Tablica Kanban została podzielona na sześć głównych etapów:

\begin{itemize}
    \item \textbf{Backlog} - przeznaczony do przechowywania wszystkich potrzebnych zadań do zrealizowania projektu. Te zadania obejmują duży zakres obowiązków, które są w następnym etapie rozdzielane na mniejsze, bardziej precyzyjne zadania, jeżeli są wybrane do realizacji w najbliższym sprincie.
    \item \textbf{Backlog sprintu} - zbierane są tu zadania do bieżącego sprintu. Dokładnie jest określany zakres wymagań, który musi spełnić, aby w pełni je zrealizować.
    \item \textbf{Wykonanie} - aktualnie wykonywane zadania. Przydzielone są one do konkretnego członka zespołu i maja ustawiony priorytet oraz liczbę ustalonych story points'ów.
    \item \textbf{Testowanie} - zbierane są tutaj wykonane zadania oczekujące na przetestowanie przez innego członka zespołu albo osobę wykonującą dane zadanie. W przypadku błędów lub niespełnienia wymagań, zadanie zostaje zwrócone do poprzedniego etapu w celu wprowadzenia poprawek.
    \item \textbf{Recenzja} - wykonane zadania oczekują na recenzję przez innego członka zespołu niż ten, który go wykonywał. Pozwala to na lepszą identyfikację błędów lub potencjalnych błędów. Podobnie jak w przypadku testowania, w przypadku niespełnienia kryteriów zadanie wraca do etapu 'Wykonanie'.
    \item \textbf{Zrobione} - przeznaczony do przechowywania wykonanych, przetestowanych i zrecenzowanych zadań.
\end{itemize}

\section{Zintegrowane środowisko programistyczne - VSC, VS, Rider (Michał Pawiłojć)}

W trakcie realizacji projektu wykorzystano różne zintegrowane środowiska programistyczne IDE, dostosowane do technologii używanych w poszczególnych komponentach aplikacji. Odpowiedni dobór narzędzi znacząco przyczynił się do efektywnej pracy zespołu, ułatwiając zarówno pisanie kodu, jak i jego debugowanie, testowanie oraz integrację z systemem kontroli wersji Git.

\subsection{Backend: Visual Studio i Rider}

Głównym środowiskiem wykorzystywanym do pracy nad backendem był \textbf{Visual Studio}, w wersji Community 2022, z uwagi na jego doskonałe wsparcie dla platformy .NET i języka C\#.

Visual Studio oferowało pełny zestaw funkcjonalności ułatwiających pracę z aplikacjami ASP.NET Core, m.in.:
\begin{itemize}
  \item integrację z systemem projektów .NET,
  \item intuicyjne debugowanie,
  \item wizualizację struktur danych,
  \item automatyczne wykrywanie problemów w kodzie (analiza statyczna),
  \item możliwość tworzenia i uruchamiania testów jednostkowych,
  \item integrację z narzędziami do konteneryzacji, np. Docker Tools.
\end{itemize}

Niektórzy członkowie zespołu korzystali również z alternatywnego IDE: \textbf{JetBrains Rider}. Był on preferowany w szczególności przez osoby przyzwyczajone do narzędzi JetBrains oraz pracujące głównie w środowisku Linux lub WSL. Rider zapewniał szybkie działanie, inteligentne podpowiedzi (ang. \textit{code completion}) i integrację z systemami bazodanowymi oraz Dockerem.

Oba środowiska dobrze współpracowały z systemem kontroli wersji Git oraz pozwalały na łatwą konfigurację środowisk uruchomieniowych (launch profiles), co było kluczowe przy testowaniu aplikacji w kontenerach.

\subsection{Frontend: Visual Studio Code}

Do pracy nad frontendem, który został zaimplementowany w technologii \textbf{React} z wykorzystaniem \textbf{TypeScript}, wykorzystywano głównie \textbf{Visual Studio Code} (VS Code). Jest to lekkie, ale bardzo elastyczne narzędzie, które doskonale sprawdza się w projektach opartych na JavaScript, dzięki bogatemu ekosystemowi rozszerzeń.

VS Code ułatwiał szybkie uruchamianie aplikacji frontendowej za pomocą skryptów \texttt{npm}, a także pozwalał na debugowanie kodu przeglądarkowego poprzez wbudowany debugger.

\subsection{Wspólne praktyki i integracja IDE z projektem}

Pomimo stosowania różnych środowisk, zespół utrzymywał spójne praktyki pracy:

\begin{itemize}
  \item \textbf{Standaryzacja formatowania kodu} --- z wykorzystaniem plików konfiguracyjnych (\texttt{.editorconfig}, \texttt{.prettierrc}),
  \item \textbf{Wspólne launch profile} --- umożliwiające uruchamianie backendu lokalnie lub w kontenerze,
  \item \textbf{Integracja z GitHubem} --- dzięki której możliwa była praca zdalna, pull requesty i code review bezpośrednio z poziomu IDE,
  \item \textbf{Debugowanie lokalne i w kontenerze} --- możliwe dzięki odpowiednim rozszerzeniom i konfiguracji,
  \item \textbf{Wsparcie dla Docker Compose} --- uruchamianie całego środowiska aplikacji z poziomu IDE.
\end{itemize}

Różnorodność wykorzystywanych środowisk nie stanowiła problemu dzięki modularnej architekturze projektu oraz konteneryzacji, która gwarantowała, że kod backendu i frontend działał w identycznym środowisku niezależnie od lokalnych ustawień IDE.

Zastosowanie odpowiednich IDE w zależności od warstwy aplikacji pozwoliło zwiększyć produktywność zespołu. Visual Studio oraz Rider znakomicie sprawdziły się przy rozwoju backendu w .NET, natomiast VS Code (ang. \textit{Visual Studio Code}) był naturalnym wyborem do pracy nad frontendem w React. Elastyczność i rozszerzalność tych środowisk umożliwiła sprawną pracę zespołową oraz efektywne wdrażanie zmian w całym cyklu życia aplikacji.

\section{Projektowanie interfejsu użytkownika - Figma (Karina Wołoszyn)}

W celu zbudowania wspólnej i klarownej wizji interfejsu użytkownika, została wykorzystana Figma, czyli aplikacja webowa pozwalająca na projektowanie prototypów UI/UX (ang. \textit{User Interface/User Experience}). Narzędzie to umożliwia pracę zespołową w czasie rzeczywistym, czyniąc komunikację i wprowadzanie szybsz, i prostsze. W projekcie \textit{UniGuesser} szczególny nacisk położono na minimalizm i prostotę zaprojektowanego interfejsu. Jednocześnie uwzględniono Księgę Identyfikacji Wizualnej Politechniki Gdańskiej, aby właściwie odwzorować kolorystykę oraz rozmieszczenie elementów zgodnie z przyjętymi standardami.
Figma posiada szeroki wachlarz funkcjonalności, które w dużym stopniu usprawniają prace zespołu. Do najważniejszych z nich należą:

\begin{itemize}
    \item \textbf{Zarządzanie warstwami i grupami} - pozwala na precyzyjną organizację elementów, co jest szczególnie ważne przy tworzeniu skomplikowanej aplikacji, aby zachować jej spójność wizualną.
    \item \textbf{Tworzenie komponentów i wariantów} - umożliwia tworzenie wielu wariantów wykorzystywanych elementów UI, co pozwala na testowanie różnych rozwiązań bez konieczności zmiany kodu.
    \item \textbf{Integracje z narzędziami deweloperskimi} - wspierają płynne przejście od fazy projektowej do wdrożenia. Jednym z takich narzędzi jest GitHub, który oferuje łatwą konfigurację, możliwość dodawania zgłoszeń oraz pull requestów, wspomagając kontrolę wersji całego projektu.
    \item \textbf{Eksport zasobów i generowanie fragmentów kodu CSS} - przyspiesza implementację zaprojektowanych widoków.
    \item \textbf{Biblioteki stylów i zasobów} - pozwala na utrzymanie spójności wizualnej przechowując zestaw fontów, kolorów, efektów i rozmiarów elementów dostępnych dla całego zespołu.
\end{itemize}

Dzięki tym zaawansowanym funkcjom Figma stała się kluczowym elementem w realizacji procesu projektowania interfejsu użytkownika \textit{UniGuesser}.